\section{The Cartan construction and an exact antiunitary reversal}
\label{sec:cq-cartan-heat}

The supplied modular--Cartan manuscript \cite{cqCartan} contains both
a sector conjugation preserving a logarithmic generator and a
standard-form involution reversing the left-minus-right generator.
The following construction retains these two equations explicitly.
It constructs the negative radial direction, its domains, and its
maps to the arithmetic heat space and moving source quotients.

\subsection{The generator equation determines the reflected parameter}

Let $\mathscr H$ be a complex Hilbert space, $K$ a self-adjoint
operator, and $E_K$ its projection-valued spectral measure. Define
\begin{equation}
\begin{split}
 T(\zeta)&=e^{-i\zeta K},\qquad \zeta=a+ib\in\mathbb C,\\
 \operatorname{Dom}T(\zeta)
 &=\left\{f\in\mathscr H:
       \int_{\mathbb R}e^{2b\lambda}\,
                   d\langle E_K(\lambda)f,f\rangle<\infty\right\}.
\end{split}
\label{eq:cq-cartan-complex-domain}
\end{equation}
This is a closed densely defined normal operator for every $b$,
including every negative $b$. Density follows because the spectral
cutoffs $E_K([-n,n])f$ belong to this domain and converge to $f$.
Closedness follows from the spectral multiplication representation.

\begin{cqproposition}
Let $J$ be an antiunitary involution and suppose the equality of
self-adjoint operators $JKJ=\epsilon K$ holds, where
$\epsilon\in\{1,-1\}$; thus it includes equality of domains.
Then
\begin{equation}
 JT(\zeta)J=T(-\epsilon\overline\zeta),\qquad
 J e^{zK}J=e^{\epsilon\overline zK}.
 \label{eq:cq-cartan-parity}
\end{equation}
Both equalities include the transported operator domains.
In particular, with $\Delta=e^{-K}$, the Tomita relation
$J\Delta J=\Delta^{-1}$ is equivalent to $JKJ=-K$ and gives
\begin{equation}
 J\Delta^{i(a+ib)}J=\Delta^{i(a-ib)},\qquad
 J e^{bK}J=e^{-bK}\quad(b\in\mathbb R).
 \label{eq:cq-Tomita-radial-reversal}
\end{equation}
\end{cqproposition}

\begin{proof}
The uniqueness of the spectral resolution gives
$JE_K(B)J=E_{\epsilon K}(B)$ for every Borel set $B$.
For a simple scalar function $\phi$, antilinearity gives
$J\phi(K)J=\overline\phi(\epsilon K)$.
Uniform approximation proves this for bounded Borel functions.
For an arbitrary finite-valued Borel function, apply the bounded
identity to $\phi\,\mathbf1_{\{|\phi|\leq n\}}$ and pass to the
spectral integral. The squared-integrability condition for
$|\phi|$ is transported by $J$, proving equality of the full
domains as well as of the values. Taking
$\phi(\lambda)=e^{-i\zeta\lambda}$ and $e^{z\lambda}$ proves
\eqref{eq:cq-cartan-parity}. Applying the same real-valued
functional calculus to $-\log$ proves the equivalence involving
$\Delta$, whose kernel is zero. Substitution of $\epsilon=-1$
proves \eqref{eq:cq-Tomita-radial-reversal}.
\end{proof}

The Cartan logarithm in the supplied construction is
\[
 \mathbf L_\ell(T)=\frac{1}{2\ell}\log(T^*T),\qquad \ell>0.
\]
For the operator in \eqref{eq:cq-cartan-complex-domain}, multiplication
of the spectral functions, including the product domain, gives
$T(a+ib)^*T(a+ib)=e^{2bK}$. Indeed the product domain is precisely
the set on which the squared multiplier $e^{2b\lambda}$ is
square-integrable; its integrability also implies that of
$e^{b\lambda}$ by Cauchy--Schwarz. Consequently,
\begin{equation}
 \mathbf L_\ell(T(a+ib))=
 \begin{cases}
 (b/\ell)K & b\ne0,\quad\text{with domain }\operatorname{Dom}K,\\
 0 & b=0,\quad\text{with domain }\mathscr H.
 \end{cases}
 \label{eq:cq-cartan-log-domain}
\end{equation}
This follows by composing the real spectral functions
$\log(e^{2b\lambda})=2b\lambda$. The zero slice has been stated
separately: the everywhere-defined $\log I$ must not be replaced
by a zero operator restricted to $\operatorname{Dom}K$.

\subsection{The original heat multiplier and a genuine modular double}

Retain the original Fourier variable $v$ and set
\begin{equation}
\begin{split}
 \mathscr H_0&=L^2(\mathbb R,dv),\qquad
 (C_0f)(v)=\overline{f(-v)},\\
 (Af)(v)&=v^2f(v),\qquad
 \operatorname{Dom}A=\{f:v^2f\in L^2\},\\
 G_t&=e^{tA/4},\\
 \operatorname{Dom}G_t
 &=\left\{f:\int_{\mathbb R}
             e^{(\operatorname{Re}t)v^2/2}|f(v)|^2\,dv<\infty\right\}.
\end{split}
\label{eq:cq-Hilbert-heat-domain}
\end{equation}
The multiplication operator $A$ is nonnegative self-adjoint,
$C_0$ is an antiunitary involution, and $C_0AC_0=A$ with the
displayed domain. Thus
$C_0G_tC_0=G_{\overline t}$, which fixes each real heat parameter.
For $\operatorname{Re}t\leq0$, $G_t$ is bounded with norm one.
For $\operatorname{Re}t>0$, it is densely defined and unbounded:
unit vectors supported in $[n,n+1]$ have image norm at least
$e^{(\operatorname{Re}t)n^2/4}$. Its domain is proper, since
$f(v)=e^{-(\operatorname{Re}t)v^2/8}$ belongs to $L^2$ and not
to that domain. Multiplication gives the exact inverse statement
\begin{equation}
 \operatorname{Ran}G_t=\operatorname{Dom}G_{-t},\qquad
 G_{-t}G_t f=f\quad(f\in\operatorname{Dom}G_t).
 \label{eq:cq-heat-unbounded-inverse}
\end{equation}
For the converse inclusion in the range statement, take
$f=G_{-t}g$ for $g\in\operatorname{Dom}G_{-t}$; then
$G_tf=g\in L^2$, proving $f\in\operatorname{Dom}G_t$.
The common invariant space
$\mathscr D=C_c^\infty(\mathbb R)$ is a graph core for each
$G_t$: truncate a vector in its weighted $L^2$ graph norm,
and approximate the truncations by smooth compact functions,
using equivalence of the weighted and unweighted norms on
each fixed compact interval.

There is no antiunitary $J_0$ on this one sector satisfying
$J_0AJ_0=-A$. Such an equality would make both $A$ and $-A$
nonnegative by preservation of real quadratic forms, forcing
$\langle Af,f\rangle=0$ for every $f\in\operatorname{Dom}A$.
A nonzero smooth function supported in $(1,2)$ contradicts this.
The exact construction retaining this original $A$ is as follows.

\begin{cqtheorem}
On $\widehat{\mathscr H}=\mathscr H_0\oplus\mathscr H_0$, define
\begin{equation}
\begin{aligned}
 \widehat K&=\begin{pmatrix}A&0\\0&-A\end{pmatrix},&
 \widehat\Delta&=\begin{pmatrix}e^{-A}&0\\0&e^A\end{pmatrix},\\
 J(f,g)&=(C_0g,C_0f),&
 \widehat G_t&=e^{t\widehat K/4}
              =\begin{pmatrix}G_t&0\\0&G_{-t}\end{pmatrix}.
\end{aligned}
\label{eq:cq-modular-heat-double}
\end{equation}
Here $\operatorname{Dom}\widehat K=\operatorname{Dom}A\oplus
\operatorname{Dom}A$,
$\operatorname{Dom}\widehat\Delta=\mathscr H_0\oplus
\operatorname{Dom}e^A$, and
$\operatorname{Dom}\widehat G_t=\operatorname{Dom}G_t\oplus
\operatorname{Dom}G_{-t}$. Then
\begin{equation}
 J^2=1,\quad J\widehat KJ=-\widehat K,\quad
 J\widehat\Delta J=\widehat\Delta^{-1},\quad
 J\widehat G_tJ=\widehat G_{-\overline t}.
 \label{eq:cq-modular-double-reversal}
\end{equation}
These are the modular data of an explicitly constructed standard
closed real subspace of $\widehat{\mathscr H}$.
\end{cqtheorem}

\begin{proof}
The block formulas and $C_0AC_0=A$ prove all four identities,
including their domains. For the last assertion put $Q=e^{-A/2}$.
This is a bounded positive injective operator, with dense range,
commuting with $C_0$. Define the closed antilinear operator
\[
 S=J\widehat\Delta^{1/2},\qquad
 \operatorname{Dom}S=\mathscr H_0\oplus\operatorname{Dom}Q^{-1},
 \qquad S(f,g)=(C_0Q^{-1}g,C_0Qf).
\]
Its domain is invariant under $S$, and direct substitution gives
$S^2(f,g)=(f,g)$. Its fixed real subspace is exactly
\begin{equation}
 V=\{(f,C_0Qf):f\in\mathscr H_0\}.
 \label{eq:cq-standard-real-heat-space}
\end{equation}
Both implications follow by substituting the second component
in $S(f,g)=(f,g)$. The graph defining $V$ is closed because
$C_0Q$ is bounded. If a vector in $V$ equals $i$ times a vector
in $V$, its first components give $f=ig$, and its second
components give $-iC_0Qg=iC_0Qg$; injectivity of $Q$ gives
$f=g=0$. Hence $V\cap iV=\{0\}$.

Moreover $V+iV=\mathscr H_0\oplus\operatorname{Ran}Q
=\operatorname{Dom}S$. To prove the nontrivial inclusion,
for $x\in\mathscr H_0$ and $y\in\operatorname{Ran}Q$ put
$b=C_0Q^{-1}y$, $f=(x+b)/2$ and $g=(x-b)/(2i)$.
Then
\[
 (f,C_0Qf)+i(g,C_0Qg)
   =(f+ig,C_0Q(f-ig))=(x,y).
\]
This domain is dense since $\operatorname{Ran}Q$ is dense.
Thus $V$ is standard: it is closed real, has zero intersection
with $iV$, and has dense complex span. On that span
$S(v+iw)=v-iw$ for $v,w\in V$, as follows from antilinearity
and $S|_V=1$. Its polar decomposition is exactly
$S=J\widehat\Delta^{1/2}$. This proves the claimed modular
real-subspace construction, including its canonical involution.
\end{proof}

For $\widehat T(\zeta)=e^{-i\zeta\widehat K}$,
the Cartan logarithm is \eqref{eq:cq-cartan-log-domain} with
$K=\widehat K$, and $\widehat G_t=\widehat T(it/4)$.
In particular the involution sends every real $t$ to $-t$.
For nonzero real $t$, the double is unbounded in one of its
two sectors; its domain remains dense. The inclusion
$I_+f=(f,0)$ retains the original multiplier exactly:
\[
 \widehat G_t I_+f=I_+G_tf
       \quad(f\in\operatorname{Dom}G_t).
\]
No Hilbert completion of the unnamed arithmetic topology is
required for this construction.

\subsection{Exact ingress and involution on the actual moving quotients}

On the original entire-function space $\cqE$ put
$(C_Eh)(s)=\overline{h(\overline s)}$. This is an antilinear
involution and preserves each previously defined growth seminorm,
since $s\mapsto\overline s$ preserves $|s|$.
For the Fourier transform already used in the test ingress,
\[
 (F\phi)(s)=\int_{\mathbb R}\phi(v)e^{-isv}\,dv,
\]
the change of variable $v\mapsto-v$ proves
$C_EF=FC_0$ on $\mathscr D$. Differentiating the integral and
summing the absolutely uniformly convergent exponential series
on the compact support proves
\begin{equation}
 U_tF=FG_t,\qquad C_EU_t=U_{\overline t}C_E.
 \label{eq:cq-cartan-Fourier-ingress}
\end{equation}
The second equality also follows directly from the convergent
entire-space series for $U_t$ and the real coefficient $-1/4$.
Thus on $\cqE\oplus\cqE$ the everywhere-defined maps
\[
 \widehat U_t=(U_t,U_{-t}),\qquad
 \mathcal J(h,k)=(C_Ek,C_Eh)
\]
satisfy $\mathcal J\widehat U_t\mathcal J
=\widehat U_{-\overline t}$ for every complex $t$. The exact
diagram is expressed by
\[
 (F\oplus F)\widehat G_t=\widehat U_t(F\oplus F),\qquad
 (F\oplus F)J=\mathcal J(F\oplus F)
       \quad\text{on }\mathscr D\oplus\mathscr D.
\]
Here $\widehat G_t$ on the test space is the restriction of
the closed Hilbert operator constructed above.

Keep the actual source topology $\tau$ and actual closed
trace image $\cqN$. Define its conjugate presentation by
$\tau^c=(C_E)_*\tau$ and $N^c=C_E\cqN$.
Since $\Xi$ is real entire, $C_E(\mathbb C[s]\Xi)
=\mathbb C[s]\Xi$, and therefore
$N^c=\overline{\mathbb C[s]\Xi}^{\,\tau^c}$.
This defines the conjugate topology explicitly without assuming
$C_E$ is continuous for the unnamed original $\tau$.

\begin{cqproposition}
For each complex $t$ form the topological quotient
\begin{equation}
\begin{split}
 \widehat N_t&=U_t\cqN\oplus U_{-t}N^c,\\
 \widehat Q_t&=
 \big((\cqE,\tau_t)\oplus(\cqE,(\tau^c)_{-t})\big)/\widehat N_t,
 \qquad \tau_t=(U_t)_*\tau,\quad
 (\tau^c)_{-t}=(U_{-t})_*\tau^c .
\end{split}
\label{eq:cq-conjugate-moving-quotient}
\end{equation}
The map $\widehat U_t$ induces a complex-linear homeomorphism
$\overline{\widehat U}_t:\widehat Q_0\to\widehat Q_t$.
The involution $\mathcal J$ induces an antilinear homeomorphism
\begin{equation}
 \overline{\mathcal J}_t:\widehat Q_t\longrightarrow
       \widehat Q_{-\overline t},\qquad
 \overline{\mathcal J}_{-\overline t}\overline{\mathcal J}_t=1,
 \qquad
 \overline{\mathcal J}_t\overline{\widehat U}_t
 =\overline{\widehat U}_{-\overline t}\overline{\mathcal J}_0 .
 \label{eq:cq-actual-quotient-antiunitary-bridge}
\end{equation}
For $\iota_t(\phi,\psi)=[F\phi,F\psi]\in\widehat Q_t$,
\begin{equation}
\begin{split}
 \ker\iota_t&=F^{-1}(U_t\cqN)\oplus F^{-1}(U_{-t}N^c),\\
 \iota_t\widehat G_t&=\overline{\widehat U}_t\iota_0,\qquad
 \overline{\mathcal J}_t\iota_t=\iota_{-\overline t}J .
\end{split}
\label{eq:cq-actual-quotient-test-diagram}
\end{equation}
The inverse images here are taken inside $\mathscr D$.
\end{cqproposition}

\begin{proof}
The direct-sum heat map transports both the specified topology
and the specified closed subspace exactly, proving its quotient
homeomorphism. Next
\[
 C_EU_{-t}N^c=U_{-\overline t}\cqN,\qquad
 C_EU_t\cqN=U_{\overline t}N^c .
\]
The same identities applied to the definitions of the transported
topologies prove that $\mathcal J$ is an antilinear homeomorphism
of the two ambient spaces in \eqref{eq:cq-conjugate-moving-quotient}.
It transports their closed relation spaces exactly, so it induces
the asserted quotient map. Applying it twice is the identity
on representatives. The conjugacy follows from
$\mathcal J\widehat U_t=\widehat U_{-\overline t}\mathcal J$.
The kernel formula is exactly the condition that both Fourier
components lie in the two relation spaces. Finally
\eqref{eq:cq-cartan-Fourier-ingress} proves both equations
in \eqref{eq:cq-actual-quotient-test-diagram} on representatives.
\end{proof}

Thus the constructed Hilbert involution is antiunitary and its
actual source-quotient image is a proved antilinear homeomorphism.
No inner product is asserted on $\widehat Q_t$.
For real $t$, this construction provides the requested negative
parameter and its exact quotient target. It does not assert
the distinct frozen inclusion $U_t\cqN\subseteq\cqN$.

\subsection{The exact viscosity sign and tensor Fourier map}

Let $\nu>0$ be the original viscosity and let $\theta\geq0$ be
physical elapsed time. On $L^2(\mathbb R^3,dq)$ define
$A_\nu=-\nu\Delta_q$ by the unitary Fourier transform
\[
 (\mathcal F_3f)(k)=(2\pi)^{-3/2}
                   \int_{\mathbb R^3}f(q)e^{-ik\cdot q}\,dq .
\]
Its domain is exactly
$\{f:|k|^2\mathcal F_3f\in L^2\}$, and its multiplier is
$\nu|k|^2$. Integration by parts verifies this differential
expression on Schwartz functions, and the spectral multiplication
domain supplies the closed realization. On the frequency tensor
product $L^2(\mathbb R)^{\otimes3}=L^2(\mathbb R^3)$ one has
\begin{equation}
 \mathcal F_3 e^{-\theta A_\nu}\mathcal F_3^{-1}
 =e^{-\nu\theta|k|^2}
 =G_{-4\nu\theta}\otimes G_{-4\nu\theta}\otimes G_{-4\nu\theta}.
 \label{eq:cq-viscosity-Fourier-intertwiner}
\end{equation}
The equality holds first on simple tensors, where multiplication
of the three factors gives the middle multiplier, and then on
all of $L^2$ by their common boundedness and tensor density.
On $\mathscr D^{\otimes3}$ the original unnormalized Fourier
transform $F^{\otimes3}$ is $(2\pi)^{3/2}\mathcal F_3$ with
the variables renamed; this scalar is retained. Each transformed
test is entire in the three coordinates, and
\eqref{eq:cq-cartan-Fourier-ingress} gives
$U_{-4\nu\theta}^{\otimes3}$ on that entire tensor test space.
This proves the sign and coefficient in the viscous bridge.
The multiplier also acts componentwise on vector fields and
preserves the divergence-free subspace: in Fourier coordinates
$k\cdot\widehat u(k)=0$ is unchanged by scalar multiplication.
Its inverse at $\theta>0$ has precisely the domain
$\{f:e^{\nu\theta|k|^2}\mathcal F_3f\in L^2\}$.
These identities identify the linear viscous operator, including
its inverse domain; they do not replace the nonlinear velocity,
pressure, or forcing equations.

\subsection{The public viscous source and its exact proof status}
\label{sec:cq-public-NS-source}

The inspected public manuscript \cite{cqNS2026} has two preserved
revisions. Its earlier 165-page PDF and its later 166-page PDF
have different recorded byte hashes. The later revision is the
one served after 19:09:35 UTC on 8 September 2026; its first
three pages were inspected directly. No equivalence of the
two complete proofs is asserted. The pinned source repository
commit is recorded separately in the bibliography and provenance.

Theorem 1.1 in that manuscript is an imported source claim:
for every $\nu>0$ it asserts existence of
$f\in C_c^\infty(\mathbb R^3\times(0,\infty);\mathbb R^3)$,
a compact spatial set $K$, and smooth $u,p$ on
$\mathbb R^3\times[0,1)$ satisfying
\[
 \partial_\theta u+(u\cdot\nabla)u-\nu\Delta u+\nabla p=f,
 \qquad \nabla\cdot u=0,\qquad u(\cdot,0)=0,
\]
with the support and norm conditions
\[
 \operatorname{supp}u(\cdot,\theta)\cup
 \operatorname{supp}p(\cdot,\theta)\subseteq K,\qquad
 \sup_{0\leq\theta<1}\|u(\cdot,\theta)\|_2<\infty,\qquad
 \limsup_{\theta\uparrow1}\|u(\cdot,\theta)\|_\infty=\infty .
\]
It also asserts nonexistence of a global smooth competitor
with the same force and initial datum and uniformly bounded
kinetic energy. These assertions have not been independently
proved in this fragment or checked here by a Lean kernel.
The identities in \eqref{eq:cq-viscosity-Fourier-intertwiner}
are independently proved above and match the exact ordinary
Laplacian and positive viscosity in this source equation.

The updated manuscript credits the preceding constructions
of Diego C\'ordoba and Luis Mart\'inez-Zoroa, and the
hypodissipative work with Zheng. Separately, the primary
statement \cite{cqABStatement} attributes announced
smooth-forcing results for incompressible porous media,
Boussinesq, and three-dimensional incompressible Euler to
Levent Alp\"oge and Tristan Buckmaster, and credits the
C\'ordoba--Mart\'inez-Zoroa programme. Those are distinct
source attributions. The operator calculations here do not
decide any private-data or appropriation allegation.
